\chapter{片段轨道与多原子分子}\label{ch:b2:fragment-orbitals}

水的路易斯结构显示两个等价的 O--H 键和两个
等价的孤电子对。然而，水受紫外线照射而电离时，
电子出现四种不同的能量，而非两种；甲烷具有四个
等价的 C--H 键，却只显示两种能量。分子中的电子并不位于成对原子之间的键内，
而是占据遍布整个分子的轨道。
本章由片段的轨道构造这些轨道，
解释水为何呈角形、乙烯为何呈平面形，以及为何碳正离子
偏好周围连有烷基。

\begin{recall}
\Cref{ch:b2:diatomic-mos}：原子轨道线性组合（LCAO）、重叠与对称性、两个相互作用的能级，
成键轨道与反键轨道。本科一年级分册：价层电子对互斥模型的构型、
杂化（作为一种描述）、共振、碳正离子及其稳定性次序（
其中已预告 C--H 键与空轨道的相互作用）。
\end{recall}

\section{片段法}

\begin{definition}[片段]\label{def:b2:fragment-orbitals:fragment}
称\emph{片段}\index{片段}为分子中被单独考虑的一组
原子（如一个原子、一对氢原子、一个 \ce{CH2} 基团），并称\emph{片段
轨道}\index{片段轨道}为其中的一条轨道。分子的
轨道由两个片段的轨道两两相互作用构造，
与双原子分子的原子轨道相同。
\end{definition}

\begin{definition}[对称性匹配组合]\label{def:b2:fragment-orbitals:symmetry-adapted}
称等价原子轨道的\emph{对称性匹配组合}\index{对称性匹配组合}
是这样一种组合：分子的每一个对称操作
（如关于镜面的反射、绕轴的旋转）作用于它时，
它或保持不变，或变为其相反数；对简并组则
变换为同一组内的组合。
\end{definition}

\begin{proposition}[两个氢原子]\label{prop:b2:fragment-orbitals:h2-pair}
对两个等价氢原子，其 1s 轨道为 $h_1$ 和 $h_2$，则
对称性匹配组合为 $h_1 + h_2$ 与 $h_1 - h_2$：前者
在任一交换两个原子的操作下保持不变，后者
变号。
\end{proposition}

\begin{proof}
交换两原子的操作使 $h_1$ 变为 $h_2$，$h_2$ 变为
$h_1$：$h_1 + h_2$ 保持不变，$h_1 - h_2$ 变为 $h_2 - h_1 = -(h_1 - h_2)$。
\end{proof}

\begin{proposition}[三个与四个氢原子]\label{prop:b2:fragment-orbitals:h4}
对位于四面体顶点的四个氢原子，其组合为一个
全对称组合 $h_1 + h_2 + h_3 + h_4$ 与一组三重简并的
组合，各有一个节面，形状与 p 轨道相似。对位于三角形顶点的三个
氢原子（如氨分子中），有一个对称组合和
一对简并组合。
\end{proposition}

\admitted

简并组合的形状及其名称（a、e、t）来自
群论，将在本科三年级分册中展开；这里的名称仅作标记。

\begin{proposition}[只有同类相互作用]\label{prop:b2:fragment-orbitals:interaction-rules}
片段轨道只与另一片段中
在分子的各对称操作下行为相同的轨道发生相互作用；
相互作用的强度随重叠增大，
随能量差增大而减小。
\end{proposition}

\begin{proof}
若两轨道在某个反射下行为不同，则一者对该平面对称，另一者
对该平面反对称，它们的重叠与
共振积分为零（见\cref{prop:b2:diatomic-mos:symmetry}）。
重叠与能量差的依赖关系与\cref{thm:b2:diatomic-mos:heteronuclear}相同。
\end{proof}

\begin{method}[片段图]\label{met:b2:fragment-orbitals:fragment-diagram}
\begin{enumerate}
  \item 将分子分成两个片段，通常一个是中心原子，另一个是
        其周围的一组原子。
  \item 构造外围原子轨道的对称性匹配组合。
  \item 将两片段的轨道按其在分子的镜面与
        轴下的行为进行分类。
  \item 使每个轨道与行为相同的轨道相互作用，能量
        越接近，作用越强；剩余轨道保持为非键轨道。
  \item 填入价电子，并读出：最高占据能级，
        成键性质与构型倾向。
\end{enumerate}
\end{method}

\section{水}

将水置于 $yz$ 平面内，$z$ 轴平分 H--O--H 键角。包含 $z$ 轴的两个
镜面为分子平面 $yz$ 和
与之垂直的平面 $xz$。相对于 $xz$ 平面，该平面交换两个
氢原子，$h_1 + h_2$ 对称，$h_1 - h_2$ 反对称。氧原子上，2s 与
$2\mathrm p_z$ 对两个平面都对称；$2\mathrm p_y$（位于
分子平面内，与 $z$ 垂直）对 $xz$ 反对称；
$2\mathrm p_x$ 对分子平面反对称，这一点
氢的任何组合都不具备。

因此 $h_1 + h_2$ 与 2s、$2\mathrm p_z$ 混合，$h_1 - h_2$ 与 $2\mathrm p_y$ 混合，且
$2\mathrm p_x$ 保持非键。共得六个轨道，其中四个被八个
价电子占据。

\begin{omfigure}
\begin{tikzpicture}[font=\small]
  % O fragment (left)
  \pic (os) at (0,-2.4) {omlevel={ud}{0.7}}; \node[left] at (-0.45,-2.4) {2s};
  \pic (op1) at (-0.45,0) {omlevel={ud}{0.4}}; \pic (op2) at (0,0) {omlevel={u}{0.4}}; \pic (op3) at (0.45,0) {omlevel={u}{0.4}};
  \node[left] at (-0.7,0) {2p};
  \node at (0,-3.2) {O};
  % H2 fragment (right)
  \pic (hp) at (6,0.9) {omlevel={u}{0.7}}; \node[right] at (6.4,0.9) {$h_1 + h_2$};
  \pic (hm) at (6,1.8) {omlevel={u}{0.7}}; \node[right] at (6.4,1.8) {$h_1 - h_2$};
  \node[anchor=north west] at (6.4,0.6) {\ce{H2} 片段};
  % MOs
  \pic (m1) at (3,-2.9) {omlevel={ud}{0.7}};  \node[omsymlabel, below=0.17cm, inner sep=0.5pt] at (m1-c) {$2a_1$};
  \pic (m2) at (3,-1.55) {omlevel={ud}{0.7}};  \node[omsymlabel, below=0.17cm, inner sep=0.5pt] at (m2-c) {$1b_2$};
  \pic (m3) at (3,-0.78) {omlevel={ud}{0.7}};  \node[omsymlabel, below=0.17cm, inner sep=0.5pt] at (m3-c) {$3a_1$};
  \pic (m4) at (3,0) {omlevel={ud}{0.7}};  \node[omsymlabel, below=0.17cm, inner sep=0.5pt] at (m4-c) {$1b_1$};
  \pic (m5) at (3,2.4) {omlevel={empty}{0.7}}; \node[omsymlabel, below=0.17cm, inner sep=0.5pt] at (m5-c) {$4a_1$};
  \pic (m6) at (3,3.1) {omlevel={empty}{0.7}}; \node[omsymlabel, below=0.17cm, inner sep=0.5pt] at (m6-c) {$2b_2$};
  \draw[omcorr, densely dashed] (os-r) -- (m1-l) (op3-r) -- (m2-l) (op3-r) -- (m3-l) (op3-r) -- (m4-l) (op3-r) -- (m5-l) (op3-r) -- (m6-l);
  \draw[omcorr, densely dashed] (hp-l) -- (m1-r) (hm-l) -- (m2-r) (hp-l) -- (m3-r) (hp-l) -- (m5-r) (hm-l) -- (m6-r);
\end{tikzpicture}

{\small 水的片段图，示意图。氢原子的对称组合与氧的
2s 和 $2\mathrm p_z$ 轨道混合（$a_1$ 轨道），
反对称组合与 $2\mathrm p_y$ 混合（$b_2$）；垂直于分子平面的 $2\mathrm p_x$ 保持非键
（$1b_1$，处于氧的 2p 能级）。四个占据的价能级，对应四个光电子谱带；
最高的 $1b_1$ 能级在
\qty{12.62}{eV} 处电离。标记
$a_1$、$b_1$、$b_2$ 为名称（来自群论，将在本科三年级分册中学习）。}
% ledger: ie:H2O
\end{omfigure}

\begin{proposition}[水为何呈弯曲形]\label{prop:b2:fragment-orbitals:bent-water}
在线性的 H--O--H 中，两个与键轴垂直的氧 p 轨道
是非键的。弯曲后，其中位于分子平面内的一个与
$h_1 + h_2$ 重叠，并与 2s 轨道混合：该占据能级下降，下降幅度大于
其他占据能级的上升幅度。价电子数为八时，弯曲的分子
更稳定。
\end{proposition}

\begin{proof}[证明]
在线性分子中（$z$ 沿 H--O--H 方向），$h_1 + h_2$ 与
$2\mathrm p_x$、$2\mathrm p_y$ 的重叠为零（二者各自相对于包含键轴的平面是反对称的，
而组合是对称的）；它们构成两对简并的
非键组合。在 $yz$ 平面内弯曲时，$2\mathrm p_y$ 部分指向
氢原子：它与 $h_1 + h_2$ 的重叠从零开始增大，其所占据的能级
随之稳定（弯曲分子中的 $3a_1$ 能级）。成键的
$1b_2$ 能级重叠略有减少，能量略有升高。净效果是能级下降，
当下降的能级被占据时下降最为明显，八个电子时正是如此。
只有四个价电子的分子，例如 \ce{BeH2}，该能级为空，
因而是线性的。
\end{proof}

\section{甲烷、氨与光电子能谱}

在甲烷中，四个氢轨道给出对称组合，它
与碳的 2s 轨道匹配；三个简并组合
与其三个 2p 轨道匹配。四个成键轨道被占据：一个较低（$a_1$），三个
简并的能量较高（$t_2$）。氨含有三个氢，氮上保留一个
占据轨道，指向远离氢原子的方向，即其孤电子对，
它是最高的占据能级。

\begin{omfigure}
\begin{tikzpicture}[font=\small]
  \pic (cs) at (0,-1.6) {omlevel={ud}{0.7}}; \node[left] at (-0.45,-1.6) {2s};
  \pic (cp1) at (-0.45,0) {omlevel={u}{0.4}}; \pic (cp2) at (0,0) {omlevel={u}{0.4}}; \pic (cp3) at (0.45,0) {omlevel={empty}{0.4}};
  \node[left] at (-0.7,0) {2p};
  \node at (0,-2.5) {C};
  \pic (ha) at (6,-0.8) {omlevel={ud}{0.7}}; \node[right] at (6.4,-0.8) {$a_1$};
  \pic (ht1) at (5.55,0.4) {omlevel={u}{0.4}}; \pic (ht2) at (6,0.4) {omlevel={u}{0.4}}; \pic (ht3) at (6.45,0.4) {omlevel={empty}{0.4}};
  \node[right] at (6.7,0.4) {$t_2$};
  \node at (6,-1.6) {\ce{H4} 片段};
  \pic (m1) at (3,-2.4) {omlevel={ud}{0.7}}; \node[omsymlabel, below=0.17cm, inner sep=0.5pt] at (m1-c) {$1a_1$};
  \pic (m2a) at (2.55,-0.9) {omlevel={ud}{0.4}}; \pic (m2b) at (3,-0.9) {omlevel={ud}{0.4}}; \pic (m2c) at (3.45,-0.9) {omlevel={ud}{0.4}};
  \node[omsymlabel, below=0.17cm, inner sep=0.5pt] at (m2b-c) {$1t_2$};
  \pic (m3) at (3,1.8) {omlevel={empty}{0.7}}; \node[omsymlabel, below=0.17cm, inner sep=0.5pt] at (m3-c) {$2a_1$};
  \pic (m4a) at (2.55,2.5) {omlevel={empty}{0.4}}; \pic (m4b) at (3,2.5) {omlevel={empty}{0.4}}; \pic (m4c) at (3.45,2.5) {omlevel={empty}{0.4}};
  \node[omsymlabel, below=0.17cm, inner sep=0.5pt] at (m4b-c) {$2t_2$};
  \draw[omcorr, densely dashed] (cs-r) -- (m1-l) (cp3-r) -- (m2a-l) (cs-r) -- (m3-l) (cp3-r) -- (m4a-l);
  \draw[omcorr, densely dashed] (ha-l) -- (m1-r) (ht1-l) -- (m2c-r) (ha-l) -- (m3-r) (ht1-l) -- (m4c-r);
\end{tikzpicture}

{\small 甲烷的片段图，示意图。四个氢的对称组合
与碳的 2s 轨道匹配，三个简并组合与其 2p
轨道匹配：得到两个占据能级，因而有两个光电子谱带，上面一个
（最先电离，位于 \qty{12.61}{eV}）容纳六个电子。}
% ledger: ie:CH4
\end{omfigure}

\begin{definition}[光电子能谱]\label{def:b2:fragment-orbitals:photoelectron}
\emph{光电子能谱}\index{光电子能谱}通过测量
已知能量的单色辐射从分子中打出的电子的动能；
谱图中的每一条谱带对应于从某一占据能级中移走一个电子，
其位置给出该能级的
电离能。
\end{definition}

\begin{proposition}[库普曼斯近似]\label{prop:b2:fragment-orbitals:koopmans}
移走占据轨道中的一个电子所需的能量，
近似等于该轨道能量的负值，$I \approx -\varepsilon$。
\end{proposition}

\admitted

该近似忽略了电离后其他电子的弛豫；
其合理性将在本科三年级分册中说明。借助它，光电子谱可以转化为轨道能级图：
水有四个价带，甲烷有两个，与片段图一致
——而与水中两个等价孤对和两个等价键的图像不同，
后者描述的是同一总电子密度，
只是描述方式不同，
但不能给出被移走电子的能量。

\section{乙烯}

乙烯可以由两个 \ce{CH2} 片段构成。每个 \ce{CH2} 除了
两个 C--H 键的轨道外，还有一个指向平面内、远离氢原子方向的轨道
（与其伙伴形成 C--C $\sigma$ 键），以及一个 p 轨道
垂直于平面。两个 p 轨道组合成一个占据的 $\pi$ 和一个
空的 $\pi^*$。

\begin{omfigure}
\begin{tikzpicture}[font=\small]
  \begin{scope}
    \pic at (0,0) {omp={0.85}{90}}; \pic at (1.2,0) {omp={0.85}{90}};
    \draw[thick] (0,0) -- (1.2,0);
    \draw[thick] (0,0) -- (-0.5,-0.45) (0,0) -- (-0.55,0.35) (1.2,0) -- (1.7,-0.45) (1.2,0) -- (1.75,0.35);
    \node at (0.6,-1.1) {$\pi$（占据）};
  \end{scope}
  \begin{scope}[xshift=3.6cm]
    \pic at (0,0) {omp={0.85}{90}}; \pic at (1.2,0) {omp={-0.85}{90}};
    \draw[thick] (0,0) -- (1.2,0);
    \draw[thick] (0,0) -- (-0.5,-0.45) (0,0) -- (-0.55,0.35) (1.2,0) -- (1.7,-0.45) (1.2,0) -- (1.75,0.35);
    \node at (0.6,-1.1) {$\pi^*$（空）};
  \end{scope}
  \begin{scope}[xshift=7.4cm]
    \pic at (0,0) {omp={0.85}{90}}; \pic at (1.2,0) {omp={0.85}{0}};
    \draw[thick] (0,0) -- (1.2,0);
    \node at (0.6,-1.1) {旋转 $90^\circ$};
    \node[font=\scriptsize, align=center] at (0.6,1.3) {零重叠};
  \end{scope}
\end{tikzpicture}

{\small 乙烯的 $\pi$ 体系示意图（分子侧视，
氢原子位于平面内）。与平面垂直的两个 p 轨道
同相（$\pi$，已填充）和异相（$\pi^*$，空）组合。将一个 \ce{CH2}
旋转 $90^\circ$ 使两个 p 轨道互相垂直：它们的重叠以及 $\pi$
键随之消失。}
\end{omfigure}

$\pi$ 键使乙烯保持平面：扭转一端会使 p 轨道彼此偏离，
其重叠随扭转角的余弦而下降，在 $90^\circ$ 时
$\pi$ 键便消失。要扭转分子，须提供
$\pi$ 键的能量，这正是烯烃的顺反异构体在
室温下不能相互转化，而绕单键的旋转是自由的。

\section{超共轭}

\begin{definition}[超共轭]\label{def:b2:fragment-orbitals:hyperconjugation}
\emph{超共轭}\index{超共轭}是指已填充的
C--H 或 C--C 键的 $\sigma$ 轨道，与相邻的空或部分填充的 p
或 $\pi^*$ 轨道（与之平行）之间的相互作用。
\end{definition}

\begin{proposition}[烷基使碳正离子稳定]\label{prop:b2:fragment-orbitals:hyperconjugation}
碳正离子被相邻碳上每一个能够
与其空 p 轨道取向一致的 C--H 或 C--C 键所稳定：此类键越多，碳
正离子越稳定，由此得出次序：叔 > 仲 > 伯 > 甲基。
\end{proposition}

\begin{proof}[论证]
碳正离子中心碳的空 p 轨道与相邻的已填充 $\sigma_{\text{C--H}}$ 轨道
构成一个两能级体系，能级差 $\Delta$ 较大，
耦合 $\beta$ 较小（仅当该键与 p 轨道大致平行时才不为零）。由
\cref{thm:b2:diatomic-mos:heteronuclear}可知，已填充的 $\sigma$ 能级
下降约 $\beta^2/\Delta$，两个电子共获得该值的两倍；空能级则
上升，且无需付出代价。每个取向合适的键各自贡献其稳定化作用；甲基
无论如何转动，总有一个 C--H 键近似取向合适。
\end{proof}

\begin{omfigure}
\begin{tikzpicture}[font=\small]
  \pic at (0,0) {omp={0.9}{90}};
  \node[circle, fill=black, inner sep=1.2pt] at (0,0) {};
  \node[font=\scriptsize, anchor=north west] at (0.08,-0.05) {C$^+$};
  \draw[thick] (0,0) -- (1.4,0);
  \node[circle, fill=black, inner sep=1.2pt] at (1.4,0) {};
  \draw[thick] (1.4,0) -- (1.75,0.85);
  \pic at (1.58,0.42) {omlobe={0.75}{68}};
  \pic at (1.58,0.42) {omlobe={0.45}{248}};
  \node[font=\scriptsize] at (1.95,1.0) {H};
  \draw[thick] (1.4,0) -- (2.0,-0.45) (1.4,0) -- (1.1,-0.6);
  \draw[<->, omProp, thick] (0.35,0.8) .. controls (0.8,1.2) and (1.2,1.2) .. (1.5,0.95);
  \node[font=\scriptsize, align=left, anchor=west] at (2.4,0.4) {已填充的 $\sigma_{\text{C--H}}$ 瓣 \\与空 p 轨道平行};
  % level scheme
  \begin{scope}[xshift=7.2cm, yshift=-0.4cm]
    \pic (p) at (0,1.4) {omlevel={empty}{0.7}}; \node[left] at (-0.4,1.4) {空 p};
    \pic (s) at (2.2,-0.8) {omlevel={ud}{0.7}}; \node[right] at (2.6,-0.8) {$\sigma_{\text{C--H}}$};
    \pic (lo) at (1.1,-1.2) {omlevel={ud}{0.7}};
    \pic (hi) at (1.1,1.7) {omlevel={empty}{0.7}};
    \draw[omcorr, densely dashed] (p-r) -- (lo-l) (p-r) -- (hi-l) (s-l) -- (lo-r) (s-l) -- (hi-r);
    \node[font=\scriptsize, anchor=north] at (1.1,-1.35) {稳定化能约为 $\approx \beta^2/\Delta$};
  \end{scope}
\end{tikzpicture}

{\small 碳正离子中超共轭的示意图。相邻碳上与空 p 轨道平行的已填充
C--H 成键轨道将部分电子密度让给空 p 轨道；已填充的能级
得到稳定，与任何两个能量不同能级之间的相互作用
一样。}
\end{omfigure}

\section{练习}

\begin{exercise}[$\star$]\label{exo:b2:fragment-orbitals:1}
写出水分子中两个氢原子 1s 轨道的两种对称性匹配组合，
并说明其中哪一个在垂直于分子且平分 H--O--H
键角的平面内改变符号。
\end{exercise}

\begin{exercise}[$\star$]\label{exo:b2:fragment-orbitals:2}
水有多少个价分子轨道，多少个价电子，
以及有多少个已填充的价轨道？
\end{exercise}

\begin{exercise}[$\star$]\label{exo:b2:fragment-orbitals:3}
为什么甲烷的光电子能谱显示出两个价带，其中一个
的强度是另一个的三倍？
\end{exercise}

\begin{exercise}[$\star$]\label{exo:b2:fragment-orbitals:4}
为什么乙烯的六个原子处于同一平面内？
\end{exercise}

\begin{exercise}[$\star\star$]\label{exo:b2:fragment-orbitals:5}
氨的第一电离能为 \qty{10.07}{eV}，甲烷的
\qty{12.61}{eV}。各自失去电子的是哪一个轨道，为何氨
更易电离？
\end{exercise}

\begin{exercise}[$\star\star$]\label{exo:b2:fragment-orbitals:6}
比较直线形与弯曲形水分子的占据能级，并解释为何价电子数为四的
分子（如 \ce{BeH2}）呈直线形。
\end{exercise}

\begin{exercise}[$\star\star$]\label{exo:b2:fragment-orbitals:7}
借助两能级的结果，解释下列碳正离子的稳定性次序：
\ce{CH3+}、\ce{CH3CH2+}、\ce{(CH3)2CH+}、\ce{(CH3)3C+}。
\end{exercise}

\begin{exercise}[$\star\star$]\label{exo:b2:fragment-orbitals:8}
绕键扭转角为 $\theta$ 的两个 p 轨道，其重叠
正比于 $\cos\theta$。若共振积分正比于
重叠，则 $\pi$ 稳定化能如何依赖于 $\theta$？它在何角度时
降为一半？
\end{exercise}

\begin{exercise}[$\star\star$]\label{exo:b2:fragment-orbitals:9}
硼烷 \ce{BH3} 为平面形，氨为三角锥形。利用中心原子与三个氢原子的
片段轨道，依据各自所含的电子解释二者的差别。
\end{exercise}

\begin{exercise}[$\star\star\star$]\label{exo:b2:fragment-orbitals:10}
\ce{H3+} 离子为正三角形。设 $\alpha$ 和 $\beta$ 分别为 1s 轨道的积分，
忽略重叠，求其三个轨道能量（对称组合的能量为
$\alpha + 2\beta$），并解释为何两个电子足以
使其成键。
\end{exercise}

\begin{exercise}[$\star\star\star$]\label{exo:b2:fragment-orbitals:11}
烯丙基体系 \ce{CH2=CH-CH2} 具有三个平行的 p 轨道。依据节点数推测
其三个 $\pi$ 轨道的形状（标出每个碳原子上的符号），并判断哪一个
是烯丙基负离子的最高占据能级。
\end{exercise}

\begin{exercise}[$\star\star\star$]\label{exo:b2:fragment-orbitals:12}
由与 O--C--O 轴垂直的 2p 轨道构造二氧化碳的 $\pi$ 体系：
共有几个 $\pi$ 轨道，哪些是成键的、非键的，
反键的，它们各容纳几个电子？
\end{exercise}

\section{周末问题：水的形状}

\begin{problem}[{周末问题——氧与 \ce{H2} 片段，线性 H--O--H 的
轨道、弯曲时移动的能级，以及利用库普曼斯近似的
光电子能谱}]
\label{pb:b2:fragment-orbitals:1}
数据：水的几何构型，$r(\ce{O-H}) = \qty{95.8}{pm}$，键角
$104.48^\circ$；第一电离能：\ce{H2O} \qty{12.62}{eV}，\ce{O}
\qty{13.62}{eV}。水位于 $yz$ 平面内，$z$ 沿角平分线方向。
% ledger: geo:H2O.rOH, geo:H2O.aHOH, ie:H2O, ie:O

\medskip
\noindent\textbf{第一部分——片段。}
\begin{enumerate}
  \item 水有多少个价电子？分别来自哪些原子？
  \item 写出氢 1s 轨道的两个对称性匹配组合。
  \item 哪些氧轨道对两个镜面都是对称的？
  \item 哪一个氧轨道仅相对于 $xz$ 平面是反对称的？
  \item 哪一个氧轨道在氢的组合中没有对应的伙伴？
  \item 共形成多少个分子轨道，其中有几个被占据？
\end{enumerate}

\medskip
\noindent\textbf{第二部分——线性 H--O--H。}
\begin{enumerate}[resume]
  \item 沿 $z$ 方向的线性分子中，哪一种组合与
        氧的 $2\mathrm p_z$ 相互作用？
  \item 哪一种与氧的 2s 相互作用？
  \item 哪些氧轨道保持为非键的？它们的简并度是多少？
  \item 用八个价电子填充各能级。
  \item 线性水是否具有顺磁性？
  \item 与氧的 2p 轨道相比，其最高占据能级
        位于何处？
\end{enumerate}

\medskip
\noindent\textbf{第三部分——弯曲。}
\begin{enumerate}[resume]
  \item 沿 $z$ 方向排列的线性分子在 $yz$ 平面内发生弯曲，哪一个
        非键轨道开始与 $h_1 + h_2$ 发生重叠？
  \item 它所在的已占据能级发生了什么变化？
  \item 哪一个成键能级略有升高？为什么？
  \item 总结：水为何是弯曲的？
  \item 将实测键角与 $90^\circ$（纯 p 成键）和 $109.5^\circ$ 作比较。
  \item 弯曲过程中，哪一个非键轨道未受影响？
\end{enumerate}

\medskip
\noindent\textbf{第四部分——光电子能谱。}
\begin{enumerate}[resume]
  \item 该能谱显示有几个价带？
  \item 由库普曼斯近似可知，第一个谱带来自哪一个轨道？
  \item 为何该谱带较窄，而成键轨道对应的谱带较宽？
  \item 比较水与氧原子的第一电离能，
        并解释其差别。
  \item 路易斯结构中的“两个等价孤电子对”变成了什么？
  \item 给出水的价电子电离谱带的数目，以及
        最低谱带的能量。
\end{enumerate}
\end{problem}
